\chapter{Gradient bounds and diffuse monotonicity}\label{sec:phase-monotonicity}
The energy inequality for arbitrary competitors in Lemma \ref{lem:phase-modica-energy} differs from the pointwise gradient estimate for entire solutions. We now prove the latter and derive the exact energy-ratio identity. The source is \cite{ChAC}, Appendix B, Problem 1, pp.~27--28. We use $W(s)=(1-s^2)^2/4$ and initially set $\varepsilon=1$.

\begin{lemma}[Uniform estimates for bounded entire solutions]\label{lem:phase-entire-estimates}
If $u\in C^2(\mathbb R^n)$, $|u|\leq1$, and $\Delta u=W'(u)$, then $u$ is smooth. For every integer $k\geq1$, $\sup_{\mathbb R^n}|D^ku|\leq C(n,k)$. Translates of any such sequence of solutions have a subsequence converging smoothly on compact sets to another solution with values in $[-1,1]$.
\end{lemma}
\begin{proof}
On each ball of radius two, both $u$ and $W'(u)$ are uniformly bounded. The interior Laplacian estimate of Theorem \ref{thm:laplace-w2p}, with $p>n$, gives a uniform $W^{2,p}$ bound on a smaller ball. Sobolev embedding gives a uniform $C^{1,\alpha}$ bound for some $\alpha>0$. The polynomial right-hand side then has controlled Hölder derivatives, and Theorem \ref{thm:interior-schauder} and its higher-order versions bootstrap to every order on smaller balls. Translation invariance makes the constants independent of the center. Arzel\`a--Ascoli and a diagonal exhaustion give the asserted compactness. For $n=1$ the same estimates follow directly from $u''=W'(u)$ and boundedness on intervals. This is the elliptic input in \cite{ChAC}, Exercise 5.6 and Problem 1(d); the underlying elliptic estimates remain the previously recorded external inputs.
\end{proof}

\begin{lemma}[A bounded function has arbitrarily small gradient somewhere]\label{lem:phase-small-gradient}
If $v\in C^2(\mathbb R^n)$ is bounded, then $\inf_{\mathbb R^n}|Dv|=0$.
\end{lemma}
\begin{proof}
Otherwise $|Dv|\geq a>0$. The vector field $X=Dv/|Dv|$ is continuously differentiable and bounded. Its integral curves exist for all forward time: bounded speed prevents escape to infinity in finite time, and local ODE continuation applies on compact sets. Along a curve $\gamma$,
$\frac d{dt}v(\gamma(t))=|Dv(\gamma(t))|\geq a$, contradicting boundedness. This proves the elementary step in \cite{ChAC}, Problem 1(a).
\end{proof}

\begin{lemma}[The differential inequality for the gradient defect]\label{lem:phase-gradient-defect-pde}
For a smooth solution of $\Delta u=W'(u)$, put $P=|Du|^2/2-W(u)$. On the open set $\{|Du|>0\}$,
\[
 \Delta P-\frac{2W'(u)}{|Du|^2}Du\cdot DP
 \geq\frac{|DP|^2}{|Du|^2}\geq0.
\]
\end{lemma}
\begin{proof}
Differentiation and the Euclidean Bochner identity give
$\Delta P=|D^2u|^2-(W'(u))^2$ and $DP=D^2u\,Du-W'(u)Du$. Hence
\[
 |D^2u|^2\geq\frac{|D^2u\,Du|^2}{|Du|^2}
 =\frac{|DP|^2+2W'(u)Du\cdot DP}{|Du|^2}+(W'(u))^2.
\]
Subtract the last square. This is \cite{ChAC}, Problem 1(b), with $P$ equal to one-half of the defect used there. The drift coefficients are locally bounded where the gradient is nonzero.
\end{proof}

\begin{theorem}[Modica's pointwise gradient estimate]\label{thm:phase-modica-gradient}
Every $C^2$ entire solution $u:\mathbb R^n\to[-1,1]$ of $\Delta u=W'(u)$ satisfies
\[
 \frac12|Du|^2\leq W(u).
\]
\end{theorem}
\begin{proof}
By Lemma \ref{lem:phase-entire-estimates}, $S=\sup P$ is finite. Suppose $S>0$, choose $x_j$ with $P(x_j)\to S$, and pass to a smooth local limit $v$ of $u(x+x_j)$. Then $P_v\leq S$ everywhere and $P_v(0)=S$. In particular $Dv(0)\ne0$.

On the connected component $O$ of $\{|Dv|>0\}$ containing zero, Lemma \ref{lem:phase-gradient-defect-pde} and the strong maximum principle make $P_v$ constant, equal to $S$. The principle applies locally to the elliptic operator with the displayed drift and no zero-order term; connectedness propagates constancy through $O$. On this component $|Dv|^2=2(W(v)+S)\geq2S$. A finite boundary point of $O$ would retain this lower bound by continuity and would still lie in the same open regular component, a contradiction. Thus $O$ is both open and closed in $\mathbb R^n$, so it is all of $\mathbb R^n$. This contradicts Lemma \ref{lem:phase-small-gradient}. Therefore $S\leq0$. This supplies the attained and unattained supremum arguments requested in \cite{ChAC}, Problem 1(c)--(d).
\end{proof}

\begin{corollary}[Nonpositive discrepancy at arbitrary diffuse scale]\label{cor:phase-modica-scaled}
If $u:\mathbb R^n\to[-1,1]$ solves $\varepsilon^2\Delta u=W'(u)$, then
\[
 \frac\varepsilon2|Du|^2-\frac{W(u)}\varepsilon\leq0.
\]
Thus the discrepancy measure of Definition \ref{def:phase-discrepancy} is nonpositive.
\end{corollary}
\begin{proof}
Apply Theorem \ref{thm:phase-modica-gradient} to $v(y)=u(\varepsilon y)$ and divide the resulting inequality by $\varepsilon$. The entire-domain hypothesis is retained by this rescaling.
\end{proof}

\begin{example}[A local solution need not satisfy Modica's bound]\label{ex:phase-local-gradient-bound}
The initial-value problem $u''=u^3-u$, $u(0)=0$, $u'(0)=2$ has a smooth solution on a neighborhood of zero, with $|u|<1$ after restricting the interval. At zero its defect is $P(0)=2-1/4=7/4>0$. Thus boundedness on a small domain and the Euler equation alone do not imply Theorem \ref{thm:phase-modica-gradient}. This also distinguishes the pointwise estimate from the energy inequality, which remains valid for this function.
\end{example}

\begin{theorem}[Rigidity in the pointwise gradient estimate]\label{thm:phase-modica-equality}
Under the hypotheses of Theorem \ref{thm:phase-modica-gradient}, equality at one point holds if and only if either $u\equiv\pm1$ or
\[
 u(x)=\tanh\left(\frac{e\cdot x-b}{\sqrt2}\right)
 \quad\text{for some }e\in\mathbb S^{n-1},\ b\in\mathbb R.
\]
\end{theorem}
\begin{proof}
By the ordered contact principle of Lemma \ref{lem:phase-ordered-contact}, a nonconstant solution has $-1<u<1$ everywhere. At an equality point $P=0$, $W(u)>0$ and hence $Du\ne0$. The strong maximum principle used above makes $P=0$ on that regular component. A finite boundary point of the component would have $Du=0$ and $P=0$, forcing $W(u)=0$, impossible for this nonconstant solution. Thus $P=0$ throughout space.

Let $q(t)=\tanh(t/\sqrt2)$ and $f=q^{-1}(u)$. Since $q'>0$ and $(q')^2=2W(q)$, the equality gives $|Df|=1$. The equation for $u=q(f)$ gives $q'\Delta f=0$, so $f$ is harmonic. Bochner's identity now gives $0=\Delta|Df|^2=2|D^2f|^2$. Therefore $f=e\cdot x-b$ with $|e|=1$. Conversely the displayed layers and constant wells satisfy equality directly. This proves \cite{ChAC}, Problem 1(e), including the constant-well alternatives.
\end{proof}

\begin{definition}[The diffuse stress-energy tensor]\label{def:phase-stress-energy}
For a smooth function on a Euclidean open set, put
\[
 e_\varepsilon=\frac\varepsilon2|Du|^2+\frac{W(u)}\varepsilon,
 \qquad T_{ij}=e_\varepsilon\delta_{ij}-\varepsilon D_i uD_j u.
\]
This tensor records the response of the energy to changes of the independent variable. We use $(\operatorname{div}T)_j=\sum_iD_iT_{ij}$ and $T:DX=\sum_{i,j}T_{ij}D_iX_j$. It provides the integration-by-parts identity behind \cite{ChAC}, Problem 1(f).
\end{definition}

\begin{proposition}[Stress-energy conservation]\label{prop:phase-stress-conservation}
For a smooth solution of $\varepsilon^2\Delta u=W'(u)$ on $U$, the tensor in Definition \ref{def:phase-stress-energy} satisfies
\[
 \operatorname{div}T=0,\qquad
 \int_U T:DX=0\quad(X\in C_c^1(U;\mathbb R^n)).
\]
\end{proposition}
\begin{proof}
Differentiate the two terms in $T$. The mixed Hessian terms cancel and leave
$(\operatorname{div}T)_j=(\varepsilon^{-1}W'(u)-\varepsilon\Delta u)D_j u=0$.
Integration by parts gives the test-field identity. Compact support removes the boundary term. The conclusion uses the Euler equation, not just an energy bound.
\end{proof}

\begin{theorem}[The exact diffuse energy-ratio identity]\label{thm:phase-energy-ratio}
Let $u$ be a smooth Allen--Cahn solution with parameter $\varepsilon>0$ on a neighborhood of $\overline{B_R(a)}$. For $0<r<R$, write $E(r)=E_\varepsilon(u;B_r(a))$. Then
\[
 \frac d{dr}\bigl(r^{1-n}E(r)\bigr)
 =r^{-n}\int_{B_r(a)}\left(\frac{W(u)}\varepsilon-\frac\varepsilon2|Du|^2\right)dx
 +\varepsilon r^{1-n}\int_{\partial B_r(a)}(\partial_\nu u)^2\,d\mathcal H^{n-1}.
\]
\end{theorem}
\begin{proof}
Translate $a$ to zero. Conservation of $T$ implies $\operatorname{div}(Tx)=\operatorname{tr}T=ne_\varepsilon-\varepsilon|Du|^2$. On the sphere, $x=r\nu$, so $(Tx)\cdot\nu=r(e_\varepsilon-\varepsilon(\partial_\nu u)^2)$. The divergence theorem and $E'(r)=\int_{\partial B_r}e_\varepsilon$ give
\[
 rE'(r)=(n-2)E(r)+2\int_{B_r}\frac{W(u)}\varepsilon
 +\varepsilon r\int_{\partial B_r}(\partial_\nu u)^2.
\]
Differentiate $r^{1-n}E(r)$ and combine terms to obtain the formula. Smoothness justifies it for every $r>0$ under consideration. This derives \cite{ChAC}, Problem 1(f), with the bulk term equal to minus the discrepancy density in Definition \ref{def:phase-discrepancy}; no extra factor two occurs.
\end{proof}

\begin{corollary}[Monotonicity and area-order energy growth]\label{cor:phase-energy-monotonicity}
For an entire solution $u:\mathbb R^n\to[-1,1]$ with fixed $\varepsilon>0$, the function $r^{1-n}E_\varepsilon(u;B_r(a))$ is nondecreasing for every center $a$. In particular, for $R\geq r>0$,
\[
 E_\varepsilon(u;B_R(a))\geq (R/r)^{n-1}E_\varepsilon(u;B_r(a)).
\]
\end{corollary}
\begin{proof}
Corollary \ref{cor:phase-modica-scaled} makes the bulk term in Theorem \ref{thm:phase-energy-ratio} nonnegative; the boundary term is a square. Integrate the derivative inequality. This is a diffuse counterpart of stationary area-ratio monotonicity, with an additional discrepancy term.
\end{proof}

\begin{corollary}[Finite total energy in dimension at least two]\label{cor:phase-finite-total-energy}
If $n\geq2$ and an entire solution $u:\mathbb R^n\to[-1,1]$ has finite total energy $\int_{\mathbb R^n}(\varepsilon|Du|^2/2+W(u)/\varepsilon)dx$ for one fixed $\varepsilon>0$, then $u\equiv+1$ or $u\equiv-1$.
\end{corollary}
\begin{proof}
If any ball had positive energy, the preceding growth bound would make the total energy infinite as $R\to\infty$, since $n-1>0$. Hence every ball has zero energy, and Proposition \ref{prop:phase-zero-energy} gives a constant well. The dimension restriction matters: the one-dimensional heteroclinics have finite energy by Theorem \ref{thm:phase-ode-classification}.
\end{proof}

\begin{example}[Energy ratios of a flat transition layer]\label{ex:phase-flat-density-ratio}
Let $n\geq2$ and $u(x)=q(x_n/\varepsilon)$, with $q(t)=\tanh(t/\sqrt2)$ and fixed $\varepsilon>0$. Then
\[
 \lim_{r\downarrow0}r^{1-n}E_\varepsilon(u;B_r)=0,
 \qquad
 \lim_{r\to\infty}r^{1-n}E_\varepsilon(u;B_r)=c_W\omega_{n-1}.
\]
The first limit follows because the smooth density is bounded near zero and the ball volume is of order $r^n$. Slicing in $x_n=t$ writes the ratio as
\[
 \omega_{n-1}\int_{-r}^{r}\varepsilon^{-1}(q'(t/\varepsilon))^2
 (1-t^2/r^2)^{(n-1)/2}\,dt.
\]
Dominated convergence and the total profile energy $c_W$ give the second limit. The ratio therefore is not constant at the diffuse scale, even for a flat layer. This computes \cite{ChAC}, Problem 1(g).
\end{example}
